% RESEARCH_DRAFT_SHA256: 70ea302bfbe00c1432285d7bb6c1d1fa37c335651c92e8afa32f9607ca70b3bd
% Layout and source-reference wrappers expanded for the web renderer.
\documentclass[11pt]{article}
\usepackage[margin=25mm]{geometry}
\usepackage{fontspec,unicode-math}
\setmainfont{Latin Modern Roman}
\setmathfont{Latin Modern Math}
\usepackage{amsmath,amssymb,mathtools,xurl,hyperref}
\hypersetup{colorlinks=true,urlcolor=blue,linkcolor=blue}
\setcounter{secnumdepth}{0}
\setlength{\parindent}{0pt}
\setlength{\parskip}{5pt}
\setlength{\emergencystretch}{3em}
\newcommand{\sref}[2]{\hyperlink{source:#1:#2}{[S#2]}}
\newcommand{\cataloguescope}[1]{\paragraph{Source record.}#1}
\newcommand{\floor}[1]{\left\lfloor#1\right\rfloor}
\begin{document}
\section*{Research-attempt provenance}
This standalone research attempt retains the complete original definition body below,
including its historical source and preparation notes. Those notes describe the
earlier catalogue preparation; the new attempt and its verification are identified
separately after that body. No original ID, catalogue status, or source statement is changed.

\section{Reliable Model Reduction for Nonlinear Parameterized Systems}\label{statement}

\subsection{Definitions and mathematical statement}

\paragraph{Shared notation.}Standard number systems are $\mathbb N,\mathbb Z,\mathbb Q,\mathbb R,\mathbb C$. Specialized notation and index ranges follow the cited source and the definitions below.


A state-space system evolves as $\dot x=f(x,u)$ with output $y=h(x,u)$. Model reduction seeks a lower-dimensional system preserving specified input-output properties. Observability distinguishes initial states from their output trajectories, while controllability concerns reachable states under allowed inputs.
\paragraph{Record scope and qualifications.}
The catalogue title is retained as metadata. Named certificates, contracts, bounds or frameworks in that title are not asserted as results of the original source extract.
The source target is broad: the system family, signal norm, parameter range, and acceptable error or complexity are not quantified in the supplied statement.

Source-record limitation: This is a broad model-reduction agenda without a quantified system/signal class, error target, or complexity target.
\paragraph{Statement recorded in the supplied source.}
Develop mathematically justified model-reduction methods for nonlinear, parameterized, coupled, and large-scale systems, including error bounds and reliable rules for choosing Krylov/Laguerre design parameters so that the reduced model gives a good approximation to the original system's response for arbitrary input signals while preserving relevant structure such as stability or passivity. \href{https://doi.org/10.4171/owr/2009/11}{[S2]}
\paragraph{Source status.}
The supplied catalogue labels this source record partially solved; that metadata was checked against the supplied snapshot on 2026-10-11. The label is preserved for provenance and is not an independent verification of present open status or mathematical correctness. The catalogue literature assessment is dated 2026-08-21.
\subsection{Short English statement}
Develop reliable nonlinear model reduction with response-error guarantees and structural preservation.
\subsection{Sources}
\begin{itemize}
\item[S1] Ulam AI and the UnsolvedMath Contributors, UnsolvedMath: A Curated Collection of Open Mathematics Problems, catalogue entry. Catalogue curation and metadata: CC BY 4.0. The source extract is retained; definitions, English summaries and qualifications are editorial additions.
\url{https://huggingface.co/datasets/ulamai/UnsolvedMath}.
\item[S2] Control Theory: On the Way to New Application Fields (2009), Oberwolfach Reports, DOI 10.4171/owr/2009/11, pp. 593--680
\url{https://doi.org/10.4171/owr/2009/11}.
\item[S3] Grune--Nesic model reduction
\url{https://doi.org/10.1137/S036301290240258X}.
\item Source-specific Oberwolfach Reports licences, including CC BY-SA 4.0 where applicable, are retained; see THIRD\_PARTY\_NOTICES.md.
\end{itemize}
\paragraph{Preparation and review.}
The original source and UnsolvedMath contributors retain attribution for the source extract. Codex (GPT-6; exact serving revision not exposed) assisted with definitions, English summaries, scope notes and formatting. No new human mathematical authorship or independently verified solution is asserted. Catalogue identity checked against the complete supplied registry and current repository. Distinct upstream records, including repeated questions and broad or damaged extracts, are retained for complete coverage of the source CS category. A repeated record is not a new independent conjecture. Complete mathematical-equivalence and current-literature review has not been performed. Preparation and final review follow contribution prompts C and D, with source gaps disclosed. Structural and rendering checks do not establish a new theorem or open status.
\label{end:statement}

\section{Literature and context check}
The primary report's model-reduction discussion, including printed pages 636--638, concerns nonlinear and parameterized systems, arbitrary input responses, and preservation of stability or passivity. It also describes an existing Laguerre-parameter proposal. This attempt chooses a restricted linear parameterized family with a common dissipative energy, proves its projection properties, and tests why stability alone cannot replace that assumption. It is not a general reduction method for the source's full agenda.
\section{Research attempt: Dissipative projection and an a posteriori response bound}

\paragraph{A parameter family whose structure survives projection}
Consider $\dot x=A(\mu(t))x+Bu$, $y=B^Tx$, where
\[
 A(\mu)=J(\mu)-R(\mu),\qquad J(\mu)^T=-J(\mu),\qquad R(\mu)\succeq\rho I,
 \quad\rho>0.
\]
Assume the displayed trajectories exist on the time interval compared; continuous bounded coefficients and locally integrable inputs suffice. Let $V\in\mathbb R^{n\times r}$ have orthonormal columns, and reduce by
\[
 \dot z=V^TA(\mu)Vz+V^TBu,\qquad \widehat y=B^TVz.
\]
The reduced skew part remains skew, and $V^TRV\succeq\rho I_r$. Thus for the full and reduced energy functions $E=\lVert x\rVert^2/2$ and $\widehat E=\lVert z\rVert^2/2$,
\[
 \dot E\leq-\rho\lVert x\rVert^2+u^Ty,
 \qquad \dot{\widehat E}\leq-\rho\lVert z\rVert^2+u^T\widehat y.
\]
These inequalities prove passivity in the input/output pairing displayed and exponential energy decay with zero input. They hold also for time-varying parameters because the energy matrix is the same identity for the whole family. Convex interpolation of such $A$ matrices preserves both algebraic inequalities: skew symmetry is linear, and a nonnegative weighted average of matrices at least $\rho I$ remains at least $\rho I$.

\paragraph{A residual bound valid for arbitrary specified inputs}
Lift the reduced state to $\widehat x=Vz$ and define its computable residual
\[
 r=\dot{\widehat x}-A(\mu)\widehat x-Bu
   =-(I-VV^T)(A(\mu)Vz+Bu).
\]
For $e=x-\widehat x$ we have $\dot e=A(\mu)e-r$. Let $U(t,s)$ be the homogeneous transition operator. Its action on any initial vector obeys
$d\lVert U(t,s)v\rVert^2/dt\leq-2\rho\lVert U(t,s)v\rVert^2$;
multiplying by $e^{2\rho t}$ and integrating gives $\lVert U(t,s)\rVert\leq e^{-\rho(t-s)}$. Variation of constants, verified by differentiating its expression, gives
\[
 e(t)=U(t,0)e(0)-\int_0^tU(t,s)r(s)\,ds.
\]
Taking norms therefore proves
\[
 \lVert y(t)-\widehat y(t)\rVert
 \leq\lVert B\rVert\left(e^{-\rho t}\lVert e(0)\rVert+
       \int_0^te^{-\rho(t-s)}\lVert r(s)\rVert\,ds\right).
 \tag{1}
\]
This is an a posteriori bound for each input and parameter trajectory; it does not predict that a chosen small subspace has a small residual. Residual evaluation can itself require full-order work, another limitation for practical large systems.

\paragraph{Two obstructions to extending this route blindly}
A stable full matrix need not yield a stable orthogonal projection. Let
\[
 A=\begin{pmatrix}-1&M\\0&-1\end{pmatrix},\qquad
 V=2^{-1/2}(1,1)^T,
\]
with $M>2$. The full solution matrix is $e^{-t}\begin{pmatrix}1&Mt\\0&1\end{pmatrix}$, which tends to zero, but $V^TAV=-1+M/2>0$ makes the reduced zero-input solution grow. Thus the common dissipativity condition, or another structure-preserving mechanism, cannot be inferred from eigenvalue stability alone.

Also, for linear systems with zero initial condition, a uniform absolute error bound over all input amplitudes is impossible unless the two input/output operators agree exactly. If one input $u$ gives a nonzero output error, the input $cu$ gives $c$ times that error. Letting $c$ increase violates any fixed finite absolute tolerance. A meaningful arbitrary-input guarantee must therefore scale with an input norm, use a bounded input class, or be a bound like (1) depending on the actual residual.

\section{Verification and adversarial audit}
The residual sign is checked by expanding the lifted reduced equation. The energy argument works for a varying parameter because no parameter-dependent Lyapunov matrix is differentiated. The counterexample is explicitly stable despite a nonnormal transient, and its one-dimensional projection is explicitly unstable. Arbitrary input signals are not silently restricted to the impulse response; their size enters the error control.
\paragraph{Audit corrections.} Required existence of compared trajectories, kept residual evaluation cost visible, and separated absolute input-independent accuracy from scale-sensitive response bounds.
\paragraph{External computation.} None; symbolic arguments and boundary checks only.
\section{Scope and remaining gaps}
The broad nonlinear/coupled agenda, automatic dimension choice, Krylov/Laguerre parameter selection, computational savings, and uniform error guarantees beyond this common-energy family remain unresolved. The displayed structure and error results are standard restricted calculations rederived here, not a new general solution.
\paragraph{Integration review.} Codex (GPT-6; exact serving revision not exposed) reviewed the written arguments and expanded layout and source-reference wrappers. Repository metadata and the production web renderer were checked using Python 3.13 and JavaScript with MathJax 3.2.2. This model-assisted review is not independent human proof verification; no PDF compilation or formal proof checking is asserted.

\section{Completion Estimate}
\noindent\textbf{COMPLETION: 0\%}
The calculations below the preserved source statement investigate a concrete route,
but establish no substantiated new progress toward the entire original target.
Known special cases and failed sufficient conditions are not credited as a solution.
\section{References}
The original source references and catalogue attribution are preserved in the
Sources section above. The following URLs were actually consulted for this attempt;
retention of an additional original reference does not claim fresh consultation of it.
\begin{itemize}
\item \url{https://ems.press/content/serial-article-files/46211} (consulted 2026-10-11). 
\end{itemize}
\section{Model and contributor disclosure}
The parent Codex GPT-6 model has reviewed the written mathematical arguments at their declared scopes. This is model-assisted review, not independent human verification.
OpenAI gpt-6-astra, requested reasoning configuration ultra, generated this research
attempt on 2026-10-11. The exact serving revision is not exposed. Codex GPT-6
performed orchestration and repository-format review; its exact serving revision
is also not exposed. The mathematical route, boundary checks and final adversarial
self-audit followed contribution prompts A then B in
\texttt{docs/contribution\_prompts.md}, after reading \texttt{CONTRIBUTING.md}.
Python 3.14.2 assembled the standalone documents; any mathematical computation is
specified in Verification. No theorem prover was used. No human mathematical
authorship, substantial human proof correction, or independent human proof checking
is asserted. Model-assisted review, source retrieval, and structural checks do not
establish correctness or novelty. Full private reasoning transcripts are unavailable
as reproducibility artifacts; the complete arguments, inputs and outcomes actually
used are stated here or in the companion batch audit.
\section{Final status}
Unresolved as a new research contribution. This is a concrete route investigation,
not a claim of new complete resolution or a current-status certification. A known
answer reproduced for context does not become a new contribution. The catalogue's
source-status label is preserved; this attempt status does not assert that a
source-labelled solved question is open.
\end{document}
